\documentclass[12pt,a4paper]{article}

\usepackage[utf8]{inputenc}
\usepackage[T1]{fontenc}
\usepackage[brazil]{babel}
\usepackage{geometry}
\usepackage{booktabs}
\usepackage{hyperref}
\usepackage{float}
\usepackage{graphicx}
\usepackage{pgfplots}
\usepackage{amsmath}
\usepackage{setspace}
\usepackage{longtable}
\usepackage{array}
\usepackage{tabularx}
\usepackage{enumitem}
\usepackage{xcolor}

\geometry{margin=2.5cm}
\pgfplotsset{compat=1.18}
\onehalfspacing
\setlength{\parindent}{1.25cm}

\title{Laborat\'orio 03: Caracteriza\c{c}\~ao da Atividade de Code Review no GitHub}
\author{Andr\'e Xavier Lazarini & Jo\~ao Pedro Guimar\~aes Ribeiro}
\date{15 de maio de 2026}

\begin{document}

\maketitle

\begin{abstract}
Este relat\'orio apresenta uma an\'alise quantitativa da atividade de \textit{code review} em \textit{pull requests} (PRs) do GitHub. A partir de um recorte dos dados coletados automaticamente via API REST, analisamos 16.646 PRs v\'alidos distribu\'idos em 200 reposit\'orios populares. O objetivo foi investigar como tamanho do PR, tempo de an\'alise, descri\c{c}\~ao e intera\c{c}\~oes se relacionam com o desfecho final da contribui\c{c}\~ao e com o esfor\c{c}o de revis\~ao. Os resultados mostram que o tempo de an\'alise e o volume de coment\'arios s\~ao os sinais mais fortes de dificuldade de integra\c{c}\~ao, enquanto o n\'umero de coment\'arios tamb\'em se destaca como o principal preditor do n\'umero de revis\~oes.
\end{abstract}

\section{Introdu\c{c}\~ao}

\subsection{Contextualiza\c{c}\~ao}
Em projetos colaborativos, o \textit{code review} funciona como um mecanismo de controle de qualidade, compartilhamento de conhecimento e redu\c{c}\~ao de risco antes da integra\c{c}\~ao do c\'odigo. No GitHub, esse processo acontece principalmente por meio de \textit{pull requests}, que concentram discuss\~oes t\'ecnicas, revis\~oes formais, pedidos de altera\c{c}\~ao e a decis\~ao final de aceitar ou rejeitar uma contribui\c{c}\~ao.

\subsection{Problema foco do experimento}
Mesmo sendo uma pr\'atica amplamente adotada, a atividade de revis\~ao n\~ao tem o mesmo comportamento para todos os PRs. Algumas contribui\c{c}\~oes s\~ao integradas rapidamente, enquanto outras acumulam coment\'arios, revis\~oes, tempo de espera e acabam fechadas sem merge. O problema investigado neste experimento \'e identificar quais caracter\'isticas observ\'aveis dos PRs est\~ao mais associadas ao sucesso da integra\c{c}\~ao e ao esfor\c{c}o de revis\~ao.

\subsection{Quest\~oes-problema ou quest\~oes de pesquisa}
\begin{itemize}[leftmargin=1.5cm]
\item \textbf{RQ01:} Qual a rela\c{c}\~ao entre o tamanho do PR e o feedback final da revis\~ao?
\item \textbf{RQ02:} Qual a rela\c{c}\~ao entre o tempo de an\'alise e o status final do PR?
\item \textbf{RQ03:} Como o tamanho da descri\c{c}\~ao do PR se relaciona com a probabilidade de merge?
\item \textbf{RQ04:} Qual a rela\c{c}\~ao entre o volume de intera\c{c}\~oes e o desfecho final do PR?
\item \textbf{RQ05:} Qual a rela\c{c}\~ao entre o tamanho do PR e o n\'umero de revis\~oes recebidas?
\item \textbf{RQ06:} PRs que demoram mais para serem analisados recebem mais revis\~oes?
\item \textbf{RQ07:} Descri\c{c}\~oes mais longas se associam a mais ou menos revis\~oes?
\item \textbf{RQ08:} Qual a correla\c{c}\~ao entre intera\c{c}\~oes sociais e o n\'umero de revis\~oes formais?
\end{itemize}

\subsection{Hip\'otese(s)}
\begin{itemize}[leftmargin=1.5cm]
\item \textbf{H1:} PRs maiores tendem a exigir mais revis\~oes e ter menor chance de merge.
\item \textbf{H2:} PRs que permanecem abertos por mais tempo t\^em maior chance de serem fechados sem merge.
\item \textbf{H3:} Descri\c{c}\~oes mais detalhadas tendem a facilitar a integra\c{c}\~ao e reduzir rodadas extras de revis\~ao.
\item \textbf{H4:} Mais coment\'arios e participantes indicam maior complexidade, mais revis\~oes e maior risco de rejei\c{c}\~ao.
\end{itemize}

\subsection{Objetivo}
\textbf{Objetivo principal:} caracterizar empiricamente a atividade de \textit{code review} no GitHub a partir de m\'etricas extra\'idas de PRs reais.

\textbf{Objetivos espec\'ificos:}
\begin{itemize}[leftmargin=1.5cm]
\item medir como as m\'etricas coletadas se distribuem entre PRs com merge e PRs fechados sem merge;
\item avaliar a rela\c{c}\~ao entre essas m\'etricas e o n\'umero de revis\~oes recebidas;
\item comparar os resultados observados com as hip\'oteses iniciais;
\item sintetizar implica\c{c}\~oes pr\'aticas e limita\c{c}\~oes do experimento.
\end{itemize}

\section{Metodologia}

\subsection{Passo a passo do experimento}
O experimento foi executado em quatro etapas:
\begin{enumerate}[leftmargin=1.5cm]
\item sele\c{c}\~ao dos reposit\'orios populares a partir da API do GitHub;
\item coleta de PRs fechados e revisados com autentica\c{c}\~ao via token;
\item limpeza dos dados para exclus\~ao de bots e filtros de validade;
\item c\'alculo de medianas, tabelas-resumo, correla\c{c}\~oes de Spearman e gr\'aficos.
\end{enumerate}

\subsection{Decis\~oes}
Algumas decis\~oes metodol\'ogicas influenciam a leitura dos resultados:
\begin{itemize}[leftmargin=1.5cm]
\item foi considerado apenas PR com pelo menos uma revis\~ao humana;
\item PRs com menos de uma hora entre cria\c{c}\~ao e fechamento/merge foram descartados;
\item contas automatizadas foram removidas da contagem de revis\~oes;
\end{itemize}

\subsection{Materiais utilizados}
\begin{itemize}[leftmargin=1.5cm]
\item API REST do GitHub;
\item script Python \texttt{scripts/collect\_github\_prs.py};
\item bibliotecas \texttt{requests}, \texttt{pandas}, \texttt{scipy}, \texttt{matplotlib} e \texttt{seaborn};
\item ambiente local com exporta\c{c}\~ao em CSV, PNG e LaTeX.
\end{itemize}

\subsection{M\'etodos utilizados}
O experimento adotou uma abordagem quantitativa, observacional e retrospectiva. Para resumir o comportamento dos PRs, utilizamos medianas por grupo. Para investigar associa\c{c}\~oes monot\^onicas entre m\'etricas e vari\'aveis-alvo, utilizamos o coeficiente de Spearman ($\rho$), adequado para distribui\c{c}\~oes assim\'etricas e presen\c{c}a de \textit{outliers}. Adotamos $p < 0{,}05$ como crit\'erio de signific\^ancia estat\'istica.

\subsection{M\'etricas e suas unidades}
\begin{table}[H]
\centering
\caption{M\'etricas utilizadas no experimento}
\begin{tabularx}{\textwidth}{@{}p{0.27\textwidth}p{0.18\textwidth}X@{}}
\toprule
M\'etrica & Unidade & Interpreta\c{c}\~ao \\ \midrule
\texttt{files\_changed} & arquivos & quantidade de arquivos alterados no PR \\
\texttt{analysis\_time\_hours} & horas & tempo entre cria\c{c}\~ao e fechamento/merge \\
\texttt{description\_length} & caracteres & tamanho do corpo textual do PR \\
\texttt{participants\_count} & participantes & total de usu\'arios \'unicos na discuss\~ao \\
\texttt{comments\_total} & coment\'arios & soma de coment\'arios de \textit{issue} e revis\~ao \\
\texttt{reviews\_count} & revis\~oes & quantidade de revis\~oes humanas registradas \\
\texttt{merged\_flag} & bin\'ario & 1 para merge e 0 para closed \\ \bottomrule
\end{tabularx}
\end{table}

\section{Visualiza\c{c}\~ao dos Resultados}

\subsubsection{Vis\~ao geral do dataset}
A Se\c{c}\~ao 4 foi atualizada diretamente a partir do arquivo \texttt{data/pull\_requests\_dataset.csv}, considerando o estado real do dataset no momento desta vers\~ao do relat\'orio. O conjunto atual possui \textbf{16.746 PRs v\'alidos}. Desse total, \textbf{12.044} PRs terminaram com merge e \textbf{4.702} foram fechados sem merge.

\begin{table}[H]
\centering
\caption{Resumo do dataset real analisado}
\begin{tabular}{@{}lrr@{}}
\toprule
Indicador & Valor & Percentual \\ \midrule
PRs v\'alidos & -- & 100,00\% \\
Reposit\'orios distintos & -- & -- \\
PRs com merge & -- & 71,92\% \\
PRs fechados sem merge & -- & 28,08\% \\
Mediana de arquivos alterados & 2 & -- \\
Mediana de tempo de an\'alise & 86,107 h & -- \\
Mediana do tamanho da descri\c{c}\~ao & 628 caracteres & -- \\
Mediana de coment\'arios & 3 & -- \\ \bottomrule
\end{tabular}
\end{table}

De forma agregada, o conjunto de dados mostra predomin\^ancia de PRs integrados, mas ainda com uma parcela relevante de contribui\c{c}\~oes encerradas sem merge. Em termos centrais, o PR mediano do dataset possui 2 arquivos alterados, permanece pouco mais de 86 horas em an\'alise, apresenta 628 caracteres de descri\c{c}\~ao e acumula 3 coment\'arios ao longo do processo.

\subsubsection{Tabelas-resumo por grupo}
\textbf{Resultados por status final do PR.} Para responder \`as quest\~oes de pesquisa relacionadas ao feedback final da revis\~ao, comparamos as medianas das m\'etricas entre PRs com \texttt{merged} e PRs com \texttt{closed}. As maiores diferen\c{c}as entre os grupos aparecem no tempo de an\'alise e no volume de intera\c{c}\~oes, enquanto o tamanho medido por arquivos alterados e o n\'umero de participantes permanecem praticamente est\'aveis.

\begin{table}[H]
\centering
\caption{Medianas das m\'etricas por status final do PR}
\begin{tabular}{@{}lrr@{}}
\toprule
M\'etrica & Merged & Closed \\ \midrule
\texttt{files\_changed} & 2 & 2 \\
\texttt{analysis\_time\_hours} & 44,312 & 1.035,999 \\
\texttt{description\_length} & 632,5 & 620 \\
\texttt{participants\_count} & 3 & 3 \\
\texttt{comments\_total} & 2 & 4 \\
\texttt{reviews\_count} & 1 & 1 \\ \bottomrule
\end{tabular}
\end{table}

\textbf{Resultados por faixa de revis\~oes.} Para a dimens\~ao de esfor\c{c}o de revis\~ao, os PRs foram agrupados em tr\^es faixas: 1 revis\~ao, 2 revis\~oes e 3 ou mais revis\~oes. Os resultados mostram crescimento monot\^onico em tempo de an\'alise, descri\c{c}\~ao, participantes e coment\'arios, indicando que PRs com mais rodadas de revis\~ao tendem tamb\'em a exigir mais coordena\c{c}\~ao e maior perman\^encia em aberto.

\begin{table}[H]
\centering
\caption{Medianas das m\'etricas por faixa de revis\~oes recebidas}
\begin{tabular}{@{}lrrr@{}}
\toprule
M\'etrica & 1 revis\~ao & 2 revis\~oes & 3+ revis\~oes \\ \midrule
\texttt{files\_changed} & 2 & 2 & 3 \\
\texttt{analysis\_time\_hours} & 52,083 & 96,923 & 211,665 \\
\texttt{description\_length} & 529 & 674 & 853 \\
\texttt{participants\_count} & 3 & 3 & 4 \\
\texttt{comments\_total} & 1 & 3 & 9 \\ \bottomrule
\end{tabular}
\end{table}

O padr\~ao mais forte no resumo descritivo aparece em \texttt{comments\_total}, cuja mediana cresce de 1 para 9 entre a primeira e a \'ultima faixa. O mesmo comportamento ocorre com o tempo de an\'alise, que quadruplica entre PRs com uma \'unica revis\~ao e PRs com tr\^es ou mais revis\~oes.

\subsubsection{Correla\c{c}\~oes estat\'isticas}
\textbf{Correla\c{c}\~oes com o status final.} As correla\c{c}\~oes de Spearman abaixo foram recalculadas diretamente a partir do CSV real. A tabela apresenta a associa\c{c}\~ao entre cada m\'etrica e a vari\'avel bin\'aria de desfecho (\texttt{merged\_flag}), considerando \texttt{1} para merge e \texttt{0} para fechamento sem merge.

\begin{table}[H]
\centering
\caption{Correla\c{c}\~oes de Spearman entre m\'etricas e status final}
\begin{tabular}{@{}lrrr@{}}
\toprule
M\'etrica & $\rho$ & $p$-valor & $n$ \\ \midrule
\texttt{files\_changed} & 0,0716 & $<$ 0,001 & 16.746 \\
\texttt{analysis\_time\_hours} & -0,4187 & $<$ 0,001 & 16.746 \\
\texttt{description\_length} & 0,0133 & 0,0849 & 16.746 \\
\texttt{participants\_count} & -0,0766 & $<$ 0,001 & 16.746 \\
\texttt{comments\_total} & -0,1480 & $<$ 0,001 & 16.746 \\ \bottomrule
\end{tabular}
\end{table}

Entre as m\'etricas analisadas, \texttt{analysis\_time\_hours} apresentou a associa\c{c}\~ao mais forte com o status final, com correla\c{c}\~ao moderada negativa. Isso indica que PRs mais demorados tendem, com maior frequ\^encia, a terminar fechados sem merge. O total de coment\'arios tamb\'em exibiu correla\c{c}\~ao negativa estatisticamente significativa, embora com intensidade menor. J\'a o tamanho da descri\c{c}\~ao n\~ao apresentou rela\c{c}\~ao estatisticamente significativa com o desfecho final.

\textbf{Correla\c{c}\~oes com o n\'umero de revis\~oes.} Para o esfor\c{c}o de revis\~ao, as correla\c{c}\~oes foram calculadas entre as m\'etricas do PR e \texttt{reviews\_count}. Nesse caso, todas as m\'etricas observadas apresentaram associa\c{c}\~oes positivas e estatisticamente significativas.

\begin{table}[H]
\centering
\caption{Correla\c{c}\~oes de Spearman entre m\'etricas e n\'umero de revis\~oes}
\begin{tabular}{@{}lrrr@{}}
\toprule
M\'etrica & $\rho$ & $p$-valor & $n$ \\ \midrule
\texttt{files\_changed} & 0,2249 & $<$ 0,001 & 16.746 \\
\texttt{analysis\_time\_hours} & 0,1839 & $<$ 0,001 & 16.746 \\
\texttt{description\_length} & 0,1128 & $<$ 0,001 & 16.746 \\
\texttt{participants\_count} & 0,3706 & $<$ 0,001 & 16.746 \\
\texttt{comments\_total} & 0,5862 & $<$ 0,001 & 16.746 \\ \bottomrule
\end{tabular}
\end{table}

O principal destaque desse bloco foi \texttt{comments\_total}, que apresentou a correla\c{c}\~ao mais alta de toda a an\'alise ($\rho = 0{,}5862$). Isso sugere que o volume de intera\c{c}\~oes sociais est\'a fortemente associado ao aumento do n\'umero de revis\~oes formais. Participantes e tempo de an\'alise tamb\'em mostraram rela\c{c}\~oes positivas relevantes, ainda que menos intensas.


\section{Discuss\~ao dos resultados}

\subsection{Confrontar quest\~oes-pesquisa}

\subsubsection{RQ01: tamanho do PR e feedback final}
No snapshot atual, a rela\c{c}\~ao entre tamanho e desfecho final aparece fraca quando o tamanho \'e medido por \texttt{files\_changed}. A correla\c{c}\~ao foi positiva, mas muito pequena ($\rho = 0{,}0698$), e as medianas de merged e closed ficaram id\^enticas em 2 arquivos. Isso sugere que, com os dados dispon\'iveis agora, o n\'umero de arquivos sozinho n\~ao separa bem PRs aceitos de PRs rejeitados.

\begin{figure}[H]
\centering
\includegraphics[width=0.82\textwidth]{figuras/rq01.png}
\caption{RQ01: tamanho do PR vs status final.}
\end{figure}

\subsubsection{RQ02: tempo de an\'alise e status final}
PRs fechados sem merge ficaram muito mais tempo abertos que PRs integrados: 1041,689 horas contra 44,669 horas em mediana. O coeficiente de Spearman moderado e negativo ($\rho = -0{,}4184$) refor\c{c}a a ideia de que demora prolongada costuma estar associada a dificuldades de integra\c{c}\~ao, retrabalho ou perda de prioridade.

\begin{figure}[H]
\centering
\includegraphics[width=0.82\textwidth]{figuras/rq02.png}
\caption{RQ02: tempo de an\'alise vs status final.}
\end{figure}

\subsubsection{RQ03: descri\c{c}\~ao do PR e chance de merge}
O tamanho da descri\c{c}\~ao praticamente n\~ao se relacionou com o status final. As medianas foram muito pr\'oximas, com leve vantagem para PRs com merge (631 contra 621 caracteres), mas a correla\c{c}\~ao foi desprez\'ivel e estatisticamente n\~ao significativa ($\rho = 0{,}0121$, $p = 0{,}1184$). Assim, neste snapshot, o tamanho da descri\c{c}\~ao n\~ao se sustenta como indicador de sucesso de integra\c{c}\~ao.

\begin{figure}[H]
\centering
\includegraphics[width=0.82\textwidth]{figuras/rq03.png}
\caption{RQ03: descri\c{c}\~ao do PR vs status final.}
\end{figure}

\subsubsection{RQ04: intera\c{c}\~oes e desfecho do PR}
O total de coment\'arios se mostrou mais informativo que o tamanho da descri\c{c}\~ao e que o n\'umero de arquivos. PRs fechados sem merge apresentaram mediana de 4 coment\'arios, contra 2 dos PRs com merge. A correla\c{c}\~ao negativa ($\rho = -0{,}1482$) ainda \'e fraca, mas aponta que intera\c{c}\~oes mais intensas tendem a acompanhar PRs com maior atrito no processo de revis\~ao.

\begin{figure}[H]
\centering
\includegraphics[width=0.82\textwidth]{figuras/rq04.png}
\caption{RQ04: intera\c{c}\~oes vs status final.}
\end{figure}

\subsubsection{RQ05: tamanho do PR e n\'umero de revis\~oes}
Quando observamos o esfor\c{c}o de revis\~ao, o tamanho medido em arquivos alterados passa a mostrar um padr\~ao mais claro. A mediana permanece em 2 arquivos para PRs com 1 ou 2 revis\~oes, mas sobe para 3 arquivos em PRs com 3 ou mais revis\~oes. A correla\c{c}\~ao ($\rho = 0{,}2254$) \'e fraca, por\'em consistente com a hip\'otese de que contribui\c{c}\~oes maiores tendem a exigir mais ciclos de valida\c{c}\~ao.

\begin{figure}[H]
\centering
\includegraphics[width=0.82\textwidth]{figuras/rq05.png}
\caption{RQ05: tamanho do PR vs n\'umero de revis\~oes.}
\end{figure}

\subsubsection{RQ06: tempo de an\'alise e n\'umero de revis\~oes}
O tempo de an\'alise cresce de forma monot\^onica conforme o n\'umero de revis\~oes aumenta: 52,749 horas para 1 revis\~ao, 97,750 para 2 revis\~oes e 212,922 para 3 ou mais revis\~oes. A correla\c{c}\~ao positiva ($\rho = 0{,}1827$) sugere que processos com mais rodadas tamb\'em tendem a ser mais longos, embora a magnitude permane\c{c}a fraca.

\begin{figure}[H]
\centering
\includegraphics[width=0.82\textwidth]{figuras/rq06.png}
\caption{RQ06: tempo de an\'alise vs n\'umero de revis\~oes.}
\end{figure}

\subsubsection{RQ07: descri\c{c}\~ao do PR e n\'umero de revis\~oes}
Descri\c{c}\~oes mais longas apareceram associadas a mais revis\~oes, e n\~ao a menos revis\~oes. A mediana cresce de 529 para 672 e depois para 851 caracteres conforme o PR recebe mais rodadas de revis\~ao. A correla\c{c}\~ao positiva ($\rho = 0{,}1127$) \'e pequena, mas enfraquece a hip\'otese de que detalhamento textual, isoladamente, reduz o esfor\c{c}o de revis\~ao.

\begin{figure}[H]
\centering
\includegraphics[width=0.82\textwidth]{figuras/rq07.png}
\caption{RQ07: descri\c{c}\~ao do PR vs n\'umero de revis\~oes.}
\end{figure}

\subsubsection{RQ08: intera\c{c}\~oes sociais e n\'umero de revis\~oes}
O total de coment\'arios cresce de 1 para 3 e depois para 9 na mediana, enquanto a correla\c{c}\~ao chega a $\rho = 0{,}5855$, a mais alta observada. Isso indica que PRs com muitas intera\c{c}\~oes tendem fortemente a concentrar mais rodadas de revis\~ao, servindo como um indicador operacional de complexidade e retrabalho.

\begin{figure}[H]
\centering
\includegraphics[width=0.82\textwidth]{figuras/rq08.png}
\caption{RQ08: intera\c{c}\~oes sociais vs n\'umero de revis\~oes.}
\end{figure}

\subsection{Insights}
\begin{itemize}[leftmargin=1.5cm]
\item tempo de an\'alise e coment\'arios foram muito mais informativos que o simples n\'umero de arquivos alterados para diferenciar PRs aceitos de PRs rejeitados;
\item o maior marcador de esfor\c{c}o de revis\~ao foi o volume de intera\c{c}\~oes, n\~ao a descri\c{c}\~ao textual;
\item descri\c{c}\~oes maiores n\~ao significaram integra\c{c}\~ao mais f\'acil; elas parecem acompanhar PRs mais complexos;
\item os resultados de tamanho ficaram limitados porque as m\'etricas de linhas alteradas n\~ao foram aproveit\'aveis neste snapshot.
\end{itemize}

\subsection{Compara\c{c}\~oes}
Em termos relativos, a diferen\c{c}a mais expressiva do bloco de status ocorreu em tempo de an\'alise: PRs fechados ficaram abertos cerca de 23 vezes mais tempo do que PRs com merge em mediana. No bloco de revis\~oes, a maior varia\c{c}\~ao ocorreu nos coment\'arios, que foram 9 vezes maiores em PRs com 3 ou mais revis\~oes do que em PRs com apenas uma revis\~ao.

\newpage

\subsection{Estat\'isticas}
Os coeficientes de Spearman indicaram:
\begin{itemize}[leftmargin=1.5cm]
\item associa\c{c}\~ao moderada negativa entre tempo de an\'alise e merge ($\rho = -0{,}4184$);
\item associa\c{c}\~ao fraca negativa entre coment\'arios e merge ($\rho = -0{,}1482$);
\item associa\c{c}\~ao moderada positiva entre coment\'arios e n\'umero de revis\~oes ($\rho = 0{,}5855$);
\item associa\c{c}\~oes fracas positivas entre arquivos alterados, tempo e descri\c{c}\~ao com o n\'umero de revis\~oes.
\end{itemize}

\section{Conclus\~ao}

Com base no snapshot analisado, o sucesso do \textit{pull request} est\'a mais associado \`a rapidez do processo e ao menor volume de intera\c{c}\~oes do que ao tamanho textual de sua descri\c{c}\~ao, o que converge com a literatura cl\'assica sobre revis\~ao de c\'odigo. A quantidade de coment\'arios e o tempo prolongado em aberto revelaram-se os principais indicadores de esfor\c{c}o de revis\~ao e de risco de fechamento sem merge, justificando a decis\~ao pr\'atica de priorizar contribui\c{c}\~oes de menor escopo conversacional e operacional. Para aprimorar essas an\'alises e evitar varia\c{c}\~oes temporais, trabalhos futuros devem focar em corrigir a coleta de m\'etricas relacionadas \`a contagem de linhas alteradas, segmentar os resultados segundo caracter\'isticas espec\'ificas dos reposit\'orios e linguagem, al\'em de investigar qualitativamente as discuss\~oes e testar modelos preditivos de risco.

\end{document}
